\begin{table}[H]
\caption{Capture rates at the extremes, 2024, and the association between creating value and
keeping it. Capture is gross operating surplus as a share of value added. The correlation
is computed against log value added across all 71 industries. Both extremes are partly
artefacts of the tax and subsidy treatment rather than of profitability, which is why the
paper rests on the association and not on the ranking. Source: authors' calculations from
the GDP-by-Industry accounts.}
\label{tab:capture}
\begin{tabular}{lr}
\toprule
Industry & Capture, per cent \\
\midrule
\multicolumn{2}{l}{\emph{Highest}} \\
Housing & 87.5 \\
Petroleum and coal products & 85.2 \\
Other real estate & 77.8 \\
Rental and leasing services and lessors of intangible assets & 77.3 \\
Farms & 76.8 \\
\addlinespace
\multicolumn{2}{l}{\emph{Lowest}} \\
Warehousing and storage & 4.0 \\
Securities, commodity contracts, and investments & 5.9 \\
Nursing and residential care facilities & 6.6 \\
Management of companies and enterprises & 8.9 \\
Social assistance & 9.6 \\
\midrule
\multicolumn{2}{l}{\emph{Association of size with capture, 71 industries}} \\
Pearson correlation, log value added with capture & $\phantom{-}$0.009 \\
Spearman rank correlation & $-$0.047 \\
\bottomrule
\end{tabular}
\end{table}
